\chapter{Vector calculus in $\mathbb{R}^3$ geometry}

\section{Euclidean structure of $\mathbb{R}^3$}

The Euclidean structure of R3 is nothing but the familiar “dot” product of vectors and is the fundamental structure characterizing \(\mathbb{R}^3\). We begin with a remark on notation. Any vector \(\mathbf{x} \in \mathbb{R}^3\) can be viewed as either a column or a row vector
\[
    \mathbf{x} = \begin{pmatrix}
        x_1 \\
        x_2 \\
        x_3
    \end{pmatrix}, \quad \mathbf{x} = \begin{pmatrix}
        x_1 & x_2 & x_3
    \end{pmatrix}.
\]

canonical orthonormal basis of \(\mathbb{R}^3\) is given by

component of the vector

matrices 3x3

\section{Orientational structure of \(\mathbb{R}^3\)}

For each triple of vectors \(\mathbf{x}, \mathbf{y}, \mathbf{z} \in \mathbb{R}^3\), we can define a matrix
\[
    (\mathbf{x}, \mathbf{y}, \mathbf{z}) = \begin{pmatrix}
        x_1 & y_1 & z_1 \\
        x_2 & y_2 & z_2 \\
        x_3 & y_3 & z_3
    \end{pmatrix}, \quad \mathbf{x} = \begin{pmatrix}
        \textbf{x}^T \\
        \textbf{y}^T \\
        \mathbf{z}^T
    \end{pmatrix}= \begin{pmatrix}
        x_1 & x_2 & x_3 \\
        y_1 & y_2 & y_3 \\
        z_1 & z_2 & z_3
    \end{pmatrix}.
\]

the reciprocally transposed matrices of R[3] having x, y, z as column and row vectors in the shown order, respectively. For fixed vectors x, y, the determinant of these matrices depends linearly on z. Therefore, there exists a vector x × y in R3 with the property that
\[
    \mathbf{x} \times \mathbf{y} \cdot \mathbf{z} = \det(\mathbf{x}, \mathbf{y}, \mathbf{z}) = \det\begin{pmatrix}
        \textbf{x}^T \\
        \textbf{y}^T \\
        \textbf{z}^T
    \end{pmatrix}.
\]

\dots

\[
    \begin{aligned}
        \mathbf{x} \times \mathbf{y} + \mathbf{y} \times \mathbf{x} = 0,                                                                                            \\
        \mathbf{x} \times (\mathbf{y} \times \mathbf{z}) + \mathbf{y} \times (\mathbf{z} \times \mathbf{x}) + \mathbf{z} \times (\mathbf{x} \times \mathbf{y}) = 0, \\
        (\mathbf{x} \times \mathbf{y}) \cdot (\mathbf{u} \times \mathbf{v}) = \dots                                                                                 \\
        \mathbf{x} \times (\mathbf{y} \times \mathbf{z}) = (\mathbf{x} \cdot \mathbf{z}) \mathbf{y} - (\mathbf{x} \cdot \mathbf{y}) \mathbf{z},                     \\
        \mathbf{e}_i \times \mathbf{e}_j = \sum_k \epsilon_{ijk} \mathbf{e}_k.
    \end{aligned}
\]

\begin{proof}
    We are going to prove the determinant identity...
\end{proof}

\dots
\section{determinant and trace identities}
\[
    \begin{aligned}
        (A \mathbf{x}) \times (A \mathbf{y}) \cdot (A \mathbf{z}) = \det(A) \mathbf{x} \times \mathbf{y} \cdot \mathbf{z}, \\
        (A \mathbf{x}) \times \mathbf{y} \cdot \mathbf{z} + \mathbf{x} \times (A \mathbf{y}) \cdot \mathbf{z} + \mathbf{x} \times \mathbf{y} \cdot (A \mathbf{z}) = \Tr(A) \mathbf{x} \times \mathbf{y} \cdot \mathbf{z},
    \end{aligned}
\]
\begin{proof}
    where the first one can be seen as
    \[
        (A \mathbf{x}) \times (A \mathbf{y}) \cdot (A \mathbf{z}) = \det\begin{pmatrix}
            A\mathbf{x} & A\mathbf{y} & A\mathbf{z}
        \end{pmatrix} = \det (A \begin{pmatrix}
                \mathbf{x} & \mathbf{y} & \mathbf{z} \end{pmatrix}) = \det(A) \det\begin{pmatrix}
            \mathbf{x} & \mathbf{y} & \mathbf{z}
        \end{pmatrix}.
    \]

    If \(A = 1 + \lambda A\) and we consider an infinitesimal transformation, then we can differentiate the first identity with the leibniz rule to get the lhs of the second identity
    \[
        \frac{\mathrm{d}}{\mathrm{d} \lambda} (((1+\lambda A) \mathbf{x}) \times ((1 + \lambda A) \mathbf{y}) \cdot ((1+\lambda A) \mathbf{z})) \bigg|_{\lambda = 0} = (A \mathbf{x}) \times \mathbf{y} \cdot \mathbf{z} + \mathbf{x} \times (A \mathbf{y}) \cdot \mathbf{z} + \mathbf{x} \times \mathbf{y} \cdot (A \mathbf{z}),
    \]
    and now we can differentiate the rhs of the first identity to get the rhs of the second identity
    \[
        \frac{\mathrm{d}}{\mathrm{d} \lambda} ((1 + \lambda \Tr(A)) \mathbf{x} \times \mathbf{y} \cdot \mathbf{z}) \bigg|_{\lambda = 0} = \Tr(A) \mathbf{x} \times \mathbf{y} \cdot \mathbf{z},
    \]
    since \(\mathbf{x} \times \mathbf{y} \cdot \mathbf{z}\) is independent of \(\lambda\) and if \(A = (\mathbf{a}_1, \mathbf{a}_2, \mathbf{a}_3)\)
    \[
        (1+\lambda A) = \begin{pmatrix}
            \mathbf{e}_1 + \lambda \mathbf{a}_1 & \mathbf{e}_2 + \lambda \mathbf{a}_2 & \mathbf{e}_3 + \lambda \mathbf{a}_3
        \end{pmatrix},
    \]
    so that differentiating the determinant of this matrix
    \[
        \det(1 + \lambda A) = (\mathbf{e}_1 + \lambda \mathbf{a}_1) \times (\mathbf{e}_2 + \lambda \mathbf{a}_2) \cdot (\mathbf{e}_3 + \lambda \mathbf{a}_3),
    \]
    with respect to \(\lambda\) at \(\lambda = 0\) gives us the sum of the diagonal entries of \(A\), which is exactly the trace of \(A\):
    \[
        \begin{aligned}
            \frac{\mathrm{d}}{\mathrm{d} \lambda} \det(1 + \lambda A) \bigg|_{\lambda = 0} & = \mathbf{a}_1 \cdot \mathbf{e}_2 \times \mathbf{e}_3 + \mathbf{e}_1 \times \mathbf{a}_2 \cdot \mathbf{e}_3 + \mathbf{e}_1 \times \mathbf{e}_2 \cdot \mathbf{a}_3 \\
                                                                                           & = \mathbf{a}_1 \cdot \mathbf{e}_2 \times \mathbf{e}_3 + \mathbf{a}_2 \cdot \mathbf{e}_3 \times \mathbf{e}_1 + \mathbf{a}_3 \cdot \mathbf{e}_1 \times \mathbf{e}_2 \\
                                                                                           & = \mathbf{a}_1 \cdot \mathbf{e}_1 + \mathbf{a}_2 \cdot \mathbf{e}_2 + \mathbf{a}_3 \cdot \mathbf{e}_3 = a_{11} + a_{22} + a_{33} = \Tr(A).
        \end{aligned}
    \]
\end{proof}

This identities are very useful when encountering the determinant and trace of a matrix in the context of Lie groups and Lie algebras, as we will see in the next chapters.

\section{Direct Product}

For two vectors \(\mathbf{x} = (x_1, x_2, x_3)\) and \(\mathbf{y} = (y_1, y_2, y_3)\) in \(\mathbb{R}^3\), we can define their direct product as a \(3\times 3\) matrix
\[
    \mathbf{x} \otimes \mathbf{y} = \begin{pmatrix}
        x_1\, y_1 & x_1\, y_2 & x_1\, y_3 \\
        x_2\, y_1 & x_2\, y_2 & x_2\, y_3 \\
        x_3\, y_1 & x_3\, y_2 & x_3\, y_3
    \end{pmatrix} \in \mathbb{R}(3) = \text{Mat}(3, \mathbb{R}).
\]
This operation is bilinear and has the following properties:
\[
    \begin{aligned}
        (\mathbf{x} \otimes \mathbf{y}) \mathbf{z} = (\mathbf{y} \cdot \mathbf{z}) \mathbf{x}, \\
        (\mathbf{x} \otimes \mathbf{y})^T = \mathbf{y} \otimes \mathbf{x},                     \\
        \Tr(\mathbf{x} \otimes \mathbf{y}) = \mathbf{x} \cdot \mathbf{y},                      \\
        \det(\mathbf{x} \otimes \mathbf{y}) = 0.
    \end{aligned}
\]
An eigenvector belong to the eigenvalue 0, so the direct product of two vectors is a singular matrix.

More identities can be derived from the definition of the direct product, for example
\[
    \begin{aligned}
        (A \mathbf{x}) \otimes (A \mathbf{y}) = A (\mathbf{x} \otimes \mathbf{y}) A^T, \\
        A = \sum_{i, j} A_{ij} \mathbf{e}_i \otimes \mathbf{e}_j.
    \end{aligned}
\]